% --- Archon physics formalization source begin ---
% archon:physics
% archon:covers PhyXMiniProblems/problem_phyx_mini_0981.lean
% archon:source-report reports/phyx_mini/problem_phyx_mini_0981.source.json

\chapter{Physics problem phyx\_mini\_0981}
\label{ch:PhyXMiniProblems_problem_phyx_mini_0981}

\paragraph{Problem source.}
\#\# Physical scenario

A metal bar with length \textbackslash{}( L \textbackslash{}), mass \textbackslash{}( m \textbackslash{}), and resistance \textbackslash{}( R \textbackslash{}) is placed on frictionless metal rails that are inclined at an angle \textbackslash{}( \textbackslash{}phi \textbackslash{}) above the horizontal. The rails have negligible resistance. A uniform magnetic field of magnitude \textbackslash{}( B \textbackslash{}) is directed downward as shown. The bar is released from rest and slides down the rails.

\#\# Question

After the terminal speed has been reached, at what rate is electrical energy being converted to thermal energy in the resistance of the bar?

\#\# Answer choices

- A: A: \textbackslash{}frac\{R\textasciicircum{}2mg\textasciicircum{}2\textbackslash{}tan\textasciicircum{}2\textbackslash{}phi\}\{L\textasciicircum{}2B\textasciicircum{}2\}

- B: B: \textbackslash{}frac\{Rm\textasciicircum{}2g\textasciicircum{}2\textbackslash{}tan\textasciicircum{}2\textbackslash{}phi\}\{L\textasciicircum{}2B\textasciicircum{}2\}

- C: C: \textbackslash{}frac\{Rmg\textbackslash{}tan\textbackslash{}phi\}\{L\textasciicircum{}2B\}

- D: D: \textbackslash{}frac\{mg\textbackslash{}tan\textbackslash{}phi\}\{RL\textasciicircum{}2B\textasciicircum{}2\}

\#\# Figure caption (auxiliary; use the image as primary evidence)

The image depicts a 3D diagram featuring an inclined plane with a rectangular block on it. 

Objects and relationships:
1. **Inclined Plane**: The plane is positioned at an angle denoted by \textbackslash{}( \textbackslash{}phi \textbackslash{}).
2. **Block**: A rectangular block of length \textbackslash{}( L \textbackslash{}) rests on the inclined plane.
3. **Dimensions**: The inclined length is labeled \textbackslash{}( a \textbackslash{}) and the base length is \textbackslash{}( b \textbackslash{}).

Vectors and labels:
1. **Magnetic Field Vectors (\textbackslash{}( \textbackslash{}vec\{B\} \textbackslash{}))**: Two blue arrows labeled \textbackslash{}( \textbackslash{}vec\{B\} \textbackslash{}) point downwards, indicating the direction of the magnetic field.
  
The diagram indicates relationships between the block, the inclined plane, and the magnetic field, useful for physics problems related to magnetism and mechanics.

\#\# Dataset metadata

Electromagnetic Induction; Physical Model Grounding Reasoning, Multi-Formula Reasoning

\paragraph{Recorded answer/context.}
B: B: \textbackslash{}frac\{Rm\textasciicircum{}2g\textasciicircum{}2\textbackslash{}tan\textasciicircum{}2\textbackslash{}phi\}\{L\textasciicircum{}2B\textasciicircum{}2\}

\paragraph{Figure/image path.}
/root/proposal\_for\_physic/science-mango/phyx\_mini\_run/phyx\_data/test\_image/981.png

\paragraph{Formalization target.}
create a compiling Lean file with sorry bodies at `PhyXMiniProblems/problem\_phyx\_mini\_0981.lean`. The Lean declarations must preserve the physical quantities, dimensions or dimensional roles, figure labels, governing-law hypotheses, and final relation expressed by this problem.
Use Mathlib/Physlib names found through LeanExplore where available. If a physics API is missing, introduce faithful local abstractions rather than scalar placeholder aliases.

\begin{theorem}[Physics formalization target]
\label{thm:physics:phyx_mini_0981:target}
\lean{PhyXMiniProblems.ProblemPhyXMini0981.terminal_bar_thermal_power_matches_answer_B}
\uses{def:physics:phyx-mini-0981:phyxminiproblems-problemphyxmini0981-lengthinmeters, def:physics:phyx-mini-0981:phyxminiproblems-problemphyxmini0981-massinkilograms, def:physics:phyx-mini-0981:phyxminiproblems-problemphyxmini0981-accelerationmagnitudeinmeterspersecondsquared, def:physics:phyx-mini-0981:phyxminiproblems-problemphyxmini0981-magneticfluxdensityinteslas, def:physics:phyx-mini-0981:phyxminiproblems-problemphyxmini0981-resistanceinohms, def:physics:phyx-mini-0981:phyxminiproblems-problemphyxmini0981-powerinwatts, def:physics:phyx-mini-0981:phyxminiproblems-problemphyxmini0981-inclinedrailbarsetup, def:physics:phyx-mini-0981:phyxminiproblems-problemphyxmini0981-matchesinclinedrailbarscenario, def:physics:phyx-mini-0981:phyxminiproblems-problemphyxmini0981-matchesprimaryinclinedrailfigure, def:physics:phyx-mini-0981:phyxminiproblems-problemphyxmini0981-hasphysicalinclinedrailparameters, def:physics:phyx-mini-0981:phyxminiproblems-problemphyxmini0981-satisfiesinclinedrailresistancecomposition, def:physics:phyx-mini-0981:phyxminiproblems-problemphyxmini0981-hasuniformappliedmagneticfield, def:physics:phyx-mini-0981:phyxminiproblems-problemphyxmini0981-satisfiesinclinedrailfieldgeometry, def:physics:phyx-mini-0981:phyxminiproblems-problemphyxmini0981-satisfiesterminalrailbarlaws}
The assigned autoformalize agent should translate this physics problem into Lean declarations in the covered file, with theorem and lemma proof bodies written as `by sorry`.
\end{theorem}
\begin{proof}
This is an autoformalization task, not a proof task. Produce statements that can later be proved by the physics prover without weakening the physical model.
\end{proof}

% --- Archon declaration topology begin ---
\section{Declaration topology}
\begin{definition}[acceleration Dimension]
  \label{def:physics:phyx-mini-0981:phyxminiproblems-problemphyxmini0981-accelerationdimension}
  \lean{PhyXMiniProblems.ProblemPhyXMini0981.accelerationDimension}
  Acceleration has physical dimension L T⁻²
\end{definition}
\begin{proof}
  This is the declared model definition of acceleration Dimension.
\end{proof}
\begin{definition}[magnetic Flux Density Dimension]
  \label{def:physics:phyx-mini-0981:phyxminiproblems-problemphyxmini0981-magneticfluxdensitydimension}
  \lean{PhyXMiniProblems.ProblemPhyXMini0981.magneticFluxDensityDimension}
  Magnetic flux density (tesla) has physical dimension M T⁻¹ C⁻¹
\end{definition}
\begin{proof}
  This is the declared model definition of magnetic Flux Density Dimension.
\end{proof}
\begin{definition}[electromotive Force Dimension]
  \label{def:physics:phyx-mini-0981:phyxminiproblems-problemphyxmini0981-electromotiveforcedimension}
  \lean{PhyXMiniProblems.ProblemPhyXMini0981.electromotiveForceDimension}
  Electromotive force (volt) has physical dimension M L² T⁻² C⁻¹
\end{definition}
\begin{proof}
  This is the declared model definition of electromotive Force Dimension.
\end{proof}
\begin{definition}[electric Current Dimension]
  \label{def:physics:phyx-mini-0981:phyxminiproblems-problemphyxmini0981-electriccurrentdimension}
  \lean{PhyXMiniProblems.ProblemPhyXMini0981.electricCurrentDimension}
  Electric current (ampere) has physical dimension C T⁻¹
\end{definition}
\begin{proof}
  This is the declared model definition of electric Current Dimension.
\end{proof}
\begin{definition}[electrical Resistance Dimension]
  \label{def:physics:phyx-mini-0981:phyxminiproblems-problemphyxmini0981-electricalresistancedimension}
  \lean{PhyXMiniProblems.ProblemPhyXMini0981.electricalResistanceDimension}
  Electrical resistance (ohm) has physical dimension M L² T⁻¹ C⁻²
\end{definition}
\begin{proof}
  This is the declared model definition of electrical Resistance Dimension.
\end{proof}
\begin{definition}[force Dimension]
  \label{def:physics:phyx-mini-0981:phyxminiproblems-problemphyxmini0981-forcedimension}
  \lean{PhyXMiniProblems.ProblemPhyXMini0981.forceDimension}
  Force (newton) has physical dimension M L T⁻²
\end{definition}
\begin{proof}
  This is the declared model definition of force Dimension.
\end{proof}
\begin{definition}[power Dimension]
  \label{def:physics:phyx-mini-0981:phyxminiproblems-problemphyxmini0981-powerdimension}
  \lean{PhyXMiniProblems.ProblemPhyXMini0981.powerDimension}
  Power (watt) has physical dimension M L² T⁻³
\end{definition}
\begin{proof}
  This is the declared model definition of power Dimension.
\end{proof}
\begin{definition}[Length Magnitude]
  \label{def:physics:phyx-mini-0981:phyxminiproblems-problemphyxmini0981-lengthmagnitude}
  \lean{PhyXMiniProblems.ProblemPhyXMini0981.LengthMagnitude}
  A nonnegative, unit-independent physical length
\end{definition}
\begin{proof}
  This is the declared model definition of Length Magnitude.
\end{proof}
\begin{definition}[Mass Magnitude]
  \label{def:physics:phyx-mini-0981:phyxminiproblems-problemphyxmini0981-massmagnitude}
  \lean{PhyXMiniProblems.ProblemPhyXMini0981.MassMagnitude}
  A nonnegative, unit-independent physical mass
\end{definition}
\begin{proof}
  This is the declared model definition of Mass Magnitude.
\end{proof}
\begin{definition}[Speed Magnitude]
  \label{def:physics:phyx-mini-0981:phyxminiproblems-problemphyxmini0981-speedmagnitude}
  \lean{PhyXMiniProblems.ProblemPhyXMini0981.SpeedMagnitude}
  A nonnegative, unit-independent speed
\end{definition}
\begin{proof}
  This is the declared model definition of Speed Magnitude.
\end{proof}
\begin{definition}[Acceleration Magnitude]
  \label{def:physics:phyx-mini-0981:phyxminiproblems-problemphyxmini0981-accelerationmagnitude}
  \lean{PhyXMiniProblems.ProblemPhyXMini0981.AccelerationMagnitude}
  \uses{def:physics:phyx-mini-0981:phyxminiproblems-problemphyxmini0981-accelerationdimension}
  A nonnegative, unit-independent acceleration magnitude
\end{definition}
\begin{proof}
  This is the declared model definition of Acceleration Magnitude.
\end{proof}
\begin{definition}[Tangential Acceleration]
  \label{def:physics:phyx-mini-0981:phyxminiproblems-problemphyxmini0981-tangentialacceleration}
  \lean{PhyXMiniProblems.ProblemPhyXMini0981.TangentialAcceleration}
  \uses{def:physics:phyx-mini-0981:phyxminiproblems-problemphyxmini0981-accelerationdimension}
  A signed tangential acceleration, positive down the rails
\end{definition}
\begin{proof}
  This is the declared model definition of Tangential Acceleration.
\end{proof}
\begin{definition}[Magnetic Flux Density Magnitude]
  \label{def:physics:phyx-mini-0981:phyxminiproblems-problemphyxmini0981-magneticfluxdensitymagnitude}
  \lean{PhyXMiniProblems.ProblemPhyXMini0981.MagneticFluxDensityMagnitude}
  \uses{def:physics:phyx-mini-0981:phyxminiproblems-problemphyxmini0981-magneticfluxdensitydimension}
  A nonnegative, unit-independent magnetic-flux-density magnitude
\end{definition}
\begin{proof}
  This is the declared model definition of Magnetic Flux Density Magnitude.
\end{proof}
\begin{definition}[Electromotive Force Magnitude]
  \label{def:physics:phyx-mini-0981:phyxminiproblems-problemphyxmini0981-electromotiveforcemagnitude}
  \lean{PhyXMiniProblems.ProblemPhyXMini0981.ElectromotiveForceMagnitude}
  \uses{def:physics:phyx-mini-0981:phyxminiproblems-problemphyxmini0981-electromotiveforcedimension}
  A nonnegative, unit-independent electromotive-force magnitude
\end{definition}
\begin{proof}
  This is the declared model definition of Electromotive Force Magnitude.
\end{proof}
\begin{definition}[Electric Current Magnitude]
  \label{def:physics:phyx-mini-0981:phyxminiproblems-problemphyxmini0981-electriccurrentmagnitude}
  \lean{PhyXMiniProblems.ProblemPhyXMini0981.ElectricCurrentMagnitude}
  \uses{def:physics:phyx-mini-0981:phyxminiproblems-problemphyxmini0981-electriccurrentdimension}
  A nonnegative, unit-independent electric-current magnitude
\end{definition}
\begin{proof}
  This is the declared model definition of Electric Current Magnitude.
\end{proof}
\begin{definition}[Resistance Magnitude]
  \label{def:physics:phyx-mini-0981:phyxminiproblems-problemphyxmini0981-resistancemagnitude}
  \lean{PhyXMiniProblems.ProblemPhyXMini0981.ResistanceMagnitude}
  \uses{def:physics:phyx-mini-0981:phyxminiproblems-problemphyxmini0981-electricalresistancedimension}
  A nonnegative, unit-independent electrical resistance
\end{definition}
\begin{proof}
  This is the declared model definition of Resistance Magnitude.
\end{proof}
\begin{definition}[Force Magnitude]
  \label{def:physics:phyx-mini-0981:phyxminiproblems-problemphyxmini0981-forcemagnitude}
  \lean{PhyXMiniProblems.ProblemPhyXMini0981.ForceMagnitude}
  \uses{def:physics:phyx-mini-0981:phyxminiproblems-problemphyxmini0981-forcedimension}
  A nonnegative, unit-independent force magnitude
\end{definition}
\begin{proof}
  This is the declared model definition of Force Magnitude.
\end{proof}
\begin{definition}[Power Magnitude]
  \label{def:physics:phyx-mini-0981:phyxminiproblems-problemphyxmini0981-powermagnitude}
  \lean{PhyXMiniProblems.ProblemPhyXMini0981.PowerMagnitude}
  \uses{def:physics:phyx-mini-0981:phyxminiproblems-problemphyxmini0981-powerdimension}
  A nonnegative, unit-independent rate of energy conversion
\end{definition}
\begin{proof}
  This is the declared model definition of Power Magnitude.
\end{proof}
\begin{definition}[length In Meters]
  \label{def:physics:phyx-mini-0981:phyxminiproblems-problemphyxmini0981-lengthinmeters}
  \lean{PhyXMiniProblems.ProblemPhyXMini0981.lengthInMeters}
  \uses{def:physics:phyx-mini-0981:phyxminiproblems-problemphyxmini0981-lengthmagnitude}
  Coherent-SI readout of length, in metres
\end{definition}
\begin{proof}
  This is the declared model definition of length In Meters.
\end{proof}
\begin{definition}[mass In Kilograms]
  \label{def:physics:phyx-mini-0981:phyxminiproblems-problemphyxmini0981-massinkilograms}
  \lean{PhyXMiniProblems.ProblemPhyXMini0981.massInKilograms}
  \uses{def:physics:phyx-mini-0981:phyxminiproblems-problemphyxmini0981-massmagnitude}
  Coherent-SI readout of mass, in kilograms
\end{definition}
\begin{proof}
  This is the declared model definition of mass In Kilograms.
\end{proof}
\begin{definition}[speed In Meters Per Second]
  \label{def:physics:phyx-mini-0981:phyxminiproblems-problemphyxmini0981-speedinmeterspersecond}
  \lean{PhyXMiniProblems.ProblemPhyXMini0981.speedInMetersPerSecond}
  \uses{def:physics:phyx-mini-0981:phyxminiproblems-problemphyxmini0981-speedmagnitude}
  Coherent-SI readout of speed, in metres per second
\end{definition}
\begin{proof}
  This is the declared model definition of speed In Meters Per Second.
\end{proof}
\begin{definition}[acceleration Magnitude In Meters Per Second Squared]
  \label{def:physics:phyx-mini-0981:phyxminiproblems-problemphyxmini0981-accelerationmagnitudeinmeterspersecondsquared}
  \lean{PhyXMiniProblems.ProblemPhyXMini0981.accelerationMagnitudeInMetersPerSecondSquared}
  \uses{def:physics:phyx-mini-0981:phyxminiproblems-problemphyxmini0981-accelerationmagnitude}
  Coherent-SI readout of an acceleration magnitude
\end{definition}
\begin{proof}
  This is the declared model definition of acceleration Magnitude In Meters Per Second Squared.
\end{proof}
\begin{definition}[acceleration In Meters Per Second Squared]
  \label{def:physics:phyx-mini-0981:phyxminiproblems-problemphyxmini0981-accelerationinmeterspersecondsquared}
  \lean{PhyXMiniProblems.ProblemPhyXMini0981.accelerationInMetersPerSecondSquared}
  \uses{def:physics:phyx-mini-0981:phyxminiproblems-problemphyxmini0981-tangentialacceleration}
  Signed coherent-SI readout of acceleration, in metres per second squared
\end{definition}
\begin{proof}
  This is the declared model definition of acceleration In Meters Per Second Squared.
\end{proof}
\begin{definition}[magnetic Flux Density In Teslas]
  \label{def:physics:phyx-mini-0981:phyxminiproblems-problemphyxmini0981-magneticfluxdensityinteslas}
  \lean{PhyXMiniProblems.ProblemPhyXMini0981.magneticFluxDensityInTeslas}
  \uses{def:physics:phyx-mini-0981:phyxminiproblems-problemphyxmini0981-magneticfluxdensitymagnitude}
  Coherent-SI readout of magnetic flux density, in teslas
\end{definition}
\begin{proof}
  This is the declared model definition of magnetic Flux Density In Teslas.
\end{proof}
\begin{definition}[electromotive Force In Volts]
  \label{def:physics:phyx-mini-0981:phyxminiproblems-problemphyxmini0981-electromotiveforceinvolts}
  \lean{PhyXMiniProblems.ProblemPhyXMini0981.electromotiveForceInVolts}
  \uses{def:physics:phyx-mini-0981:phyxminiproblems-problemphyxmini0981-electromotiveforcemagnitude}
  Coherent-SI readout of electromotive force, in volts
\end{definition}
\begin{proof}
  This is the declared model definition of electromotive Force In Volts.
\end{proof}
\begin{definition}[electric Current In Amperes]
  \label{def:physics:phyx-mini-0981:phyxminiproblems-problemphyxmini0981-electriccurrentinamperes}
  \lean{PhyXMiniProblems.ProblemPhyXMini0981.electricCurrentInAmperes}
  \uses{def:physics:phyx-mini-0981:phyxminiproblems-problemphyxmini0981-electriccurrentmagnitude}
  Coherent-SI readout of current, in amperes
\end{definition}
\begin{proof}
  This is the declared model definition of electric Current In Amperes.
\end{proof}
\begin{definition}[resistance In Ohms]
  \label{def:physics:phyx-mini-0981:phyxminiproblems-problemphyxmini0981-resistanceinohms}
  \lean{PhyXMiniProblems.ProblemPhyXMini0981.resistanceInOhms}
  \uses{def:physics:phyx-mini-0981:phyxminiproblems-problemphyxmini0981-resistancemagnitude}
  Coherent-SI readout of resistance, in ohms
\end{definition}
\begin{proof}
  This is the declared model definition of resistance In Ohms.
\end{proof}
\begin{definition}[force In Newtons]
  \label{def:physics:phyx-mini-0981:phyxminiproblems-problemphyxmini0981-forceinnewtons}
  \lean{PhyXMiniProblems.ProblemPhyXMini0981.forceInNewtons}
  \uses{def:physics:phyx-mini-0981:phyxminiproblems-problemphyxmini0981-forcemagnitude}
  Coherent-SI readout of force, in newtons
\end{definition}
\begin{proof}
  This is the declared model definition of force In Newtons.
\end{proof}
\begin{definition}[power In Watts]
  \label{def:physics:phyx-mini-0981:phyxminiproblems-problemphyxmini0981-powerinwatts}
  \lean{PhyXMiniProblems.ProblemPhyXMini0981.powerInWatts}
  \uses{def:physics:phyx-mini-0981:phyxminiproblems-problemphyxmini0981-powermagnitude}
  Coherent-SI readout of power, in watts
\end{definition}
\begin{proof}
  This is the declared model definition of power In Watts.
\end{proof}
\begin{definition}[Rail]
  \label{def:physics:phyx-mini-0981:phyxminiproblems-problemphyxmini0981-rail}
  \lean{PhyXMiniProblems.ProblemPhyXMini0981.Rail}
  The two frictionless rails visible in image 981.png
\end{definition}
\begin{proof}
  This is the declared model definition of Rail.
\end{proof}
\begin{definition}[Bar Endpoint]
  \label{def:physics:phyx-mini-0981:phyxminiproblems-problemphyxmini0981-barendpoint}
  \lean{PhyXMiniProblems.ProblemPhyXMini0981.BarEndpoint}
  The printed endpoint names on the two ends of the bar
\end{definition}
\begin{proof}
  This is the declared model definition of Bar Endpoint.
\end{proof}
\begin{definition}[Field Arrow]
  \label{def:physics:phyx-mini-0981:phyxminiproblems-problemphyxmini0981-fieldarrow}
  \lean{PhyXMiniProblems.ProblemPhyXMini0981.FieldArrow}
  The two blue magnetic-field arrows drawn in the primary raster
\end{definition}
\begin{proof}
  This is the declared model definition of Field Arrow.
\end{proof}
\begin{definition}[Spatial Direction]
  \label{def:physics:phyx-mini-0981:phyxminiproblems-problemphyxmini0981-spatialdirection}
  \lean{PhyXMiniProblems.ProblemPhyXMini0981.SpatialDirection}
  Qualitative directions needed to transcribe the inclined-rail geometry
\end{definition}
\begin{proof}
  This is the declared model definition of Spatial Direction.
\end{proof}
\begin{definition}[expected Rail For Endpoint]
  \label{def:physics:phyx-mini-0981:phyxminiproblems-problemphyxmini0981-expectedrailforendpoint}
  \lean{PhyXMiniProblems.ProblemPhyXMini0981.expectedRailForEndpoint}
  \uses{def:physics:phyx-mini-0981:phyxminiproblems-problemphyxmini0981-rail, def:physics:phyx-mini-0981:phyxminiproblems-problemphyxmini0981-barendpoint}
  The rail on which each labelled end of the bar lies
\end{definition}
\begin{proof}
  This is the declared model definition of expected Rail For Endpoint.
\end{proof}
\begin{definition}[Inclined Rail Figure]
  \label{def:physics:phyx-mini-0981:phyxminiproblems-problemphyxmini0981-inclinedrailfigure}
  \lean{PhyXMiniProblems.ProblemPhyXMini0981.InclinedRailFigure}
  \uses{def:physics:phyx-mini-0981:phyxminiproblems-problemphyxmini0981-rail, def:physics:phyx-mini-0981:phyxminiproblems-problemphyxmini0981-barendpoint, def:physics:phyx-mini-0981:phyxminiproblems-problemphyxmini0981-fieldarrow, def:physics:phyx-mini-0981:phyxminiproblems-problemphyxmini0981-spatialdirection}
  Literal content of the supplied raster. The symbols a and b label bar endpoints; they are not additional scalar lengths
\end{definition}
\begin{proof}
  This is the declared model definition of Inclined Rail Figure.
\end{proof}
\begin{definition}[Inclined Rail Bar Setup]
  \label{def:physics:phyx-mini-0981:phyxminiproblems-problemphyxmini0981-inclinedrailbarsetup}
  \lean{PhyXMiniProblems.ProblemPhyXMini0981.InclinedRailBarSetup}
  \uses{def:physics:phyx-mini-0981:phyxminiproblems-problemphyxmini0981-lengthmagnitude, def:physics:phyx-mini-0981:phyxminiproblems-problemphyxmini0981-massmagnitude, def:physics:phyx-mini-0981:phyxminiproblems-problemphyxmini0981-speedmagnitude, def:physics:phyx-mini-0981:phyxminiproblems-problemphyxmini0981-accelerationmagnitude, def:physics:phyx-mini-0981:phyxminiproblems-problemphyxmini0981-tangentialacceleration, def:physics:phyx-mini-0981:phyxminiproblems-problemphyxmini0981-magneticfluxdensitymagnitude, def:physics:phyx-mini-0981:phyxminiproblems-problemphyxmini0981-electromotiveforcemagnitude, def:physics:phyx-mini-0981:phyxminiproblems-problemphyxmini0981-electriccurrentmagnitude, def:physics:phyx-mini-0981:phyxminiproblems-problemphyxmini0981-resistancemagnitude, def:physics:phyx-mini-0981:phyxminiproblems-problemphyxmini0981-forcemagnitude, def:physics:phyx-mini-0981:phyxminiproblems-problemphyxmini0981-powermagnitude, def:physics:phyx-mini-0981:phyxminiproblems-problemphyxmini0981-inclinedrailfigure}
  Independent terminal observables and physical parameters of the setup
\end{definition}
\begin{proof}
  This is the declared model definition of Inclined Rail Bar Setup.
\end{proof}
\begin{definition}[Matches Inclined Rail Bar Scenario]
  \label{def:physics:phyx-mini-0981:phyxminiproblems-problemphyxmini0981-matchesinclinedrailbarscenario}
  \lean{PhyXMiniProblems.ProblemPhyXMini0981.MatchesInclinedRailBarScenario}
  \uses{def:physics:phyx-mini-0981:phyxminiproblems-problemphyxmini0981-speedinmeterspersecond, def:physics:phyx-mini-0981:phyxminiproblems-problemphyxmini0981-resistanceinohms, def:physics:phyx-mini-0981:phyxminiproblems-problemphyxmini0981-inclinedrailbarsetup}
  The apparatus properties stated in the prose, including release from rest
\end{definition}
\begin{proof}
  This is the declared model definition of Matches Inclined Rail Bar Scenario.
\end{proof}
\begin{definition}[Matches Primary Inclined Rail Figure]
  \label{def:physics:phyx-mini-0981:phyxminiproblems-problemphyxmini0981-matchesprimaryinclinedrailfigure}
  \lean{PhyXMiniProblems.ProblemPhyXMini0981.MatchesPrimaryInclinedRailFigure}
  \uses{def:physics:phyx-mini-0981:phyxminiproblems-problemphyxmini0981-expectedrailforendpoint, def:physics:phyx-mini-0981:phyxminiproblems-problemphyxmini0981-inclinedrailbarsetup}
  Primary-raster transcription for endpoints, length, incline, and field
\end{definition}
\begin{proof}
  This is the declared model definition of Matches Primary Inclined Rail Figure.
\end{proof}
\begin{definition}[Has Physical Inclined Rail Parameters]
  \label{def:physics:phyx-mini-0981:phyxminiproblems-problemphyxmini0981-hasphysicalinclinedrailparameters}
  \lean{PhyXMiniProblems.ProblemPhyXMini0981.HasPhysicalInclinedRailParameters}
  \uses{def:physics:phyx-mini-0981:phyxminiproblems-problemphyxmini0981-lengthinmeters, def:physics:phyx-mini-0981:phyxminiproblems-problemphyxmini0981-massinkilograms, def:physics:phyx-mini-0981:phyxminiproblems-problemphyxmini0981-speedinmeterspersecond, def:physics:phyx-mini-0981:phyxminiproblems-problemphyxmini0981-accelerationmagnitudeinmeterspersecondsquared, def:physics:phyx-mini-0981:phyxminiproblems-problemphyxmini0981-magneticfluxdensityinteslas, def:physics:phyx-mini-0981:phyxminiproblems-problemphyxmini0981-resistanceinohms, def:physics:phyx-mini-0981:phyxminiproblems-problemphyxmini0981-inclinedrailbarsetup}
  Positivity and nondegeneracy assumptions for the physical apparatus
\end{definition}
\begin{proof}
  This is the declared model definition of Has Physical Inclined Rail Parameters.
\end{proof}
\begin{definition}[Satisfies Inclined Rail Resistance Composition]
  \label{def:physics:phyx-mini-0981:phyxminiproblems-problemphyxmini0981-satisfiesinclinedrailresistancecomposition}
  \lean{PhyXMiniProblems.ProblemPhyXMini0981.SatisfiesInclinedRailResistanceComposition}
  \uses{def:physics:phyx-mini-0981:phyxminiproblems-problemphyxmini0981-resistanceinohms, def:physics:phyx-mini-0981:phyxminiproblems-problemphyxmini0981-inclinedrailbarsetup}
  The rails and bar resistances compose the resistance of the closed circuit
\end{definition}
\begin{proof}
  This is the declared model definition of Satisfies Inclined Rail Resistance Composition.
\end{proof}
\begin{definition}[Has Uniform Applied Magnetic Field]
  \label{def:physics:phyx-mini-0981:phyxminiproblems-problemphyxmini0981-hasuniformappliedmagneticfield}
  \lean{PhyXMiniProblems.ProblemPhyXMini0981.HasUniformAppliedMagneticField}
  \uses{def:physics:phyx-mini-0981:phyxminiproblems-problemphyxmini0981-magneticfluxdensityinteslas, def:physics:phyx-mini-0981:phyxminiproblems-problemphyxmini0981-inclinedrailbarsetup}
  The Physlib magnetic vector field is constant throughout the marked region, and its norm is calibrated by the independent dimensionful applied magnitude B
\end{definition}
\begin{proof}
  This is the declared model definition of Has Uniform Applied Magnetic Field.
\end{proof}
\begin{definition}[Satisfies Inclined Rail Field Geometry]
  \label{def:physics:phyx-mini-0981:phyxminiproblems-problemphyxmini0981-satisfiesinclinedrailfieldgeometry}
  \lean{PhyXMiniProblems.ProblemPhyXMini0981.SatisfiesInclinedRailFieldGeometry}
  \uses{def:physics:phyx-mini-0981:phyxminiproblems-problemphyxmini0981-magneticfluxdensityinteslas, def:physics:phyx-mini-0981:phyxminiproblems-problemphyxmini0981-inclinedrailbarsetup}
  The vertically downward field contributes its plane-normal component B cos φ to both motional induction and magnetic drag
\end{definition}
\begin{proof}
  This is the declared model definition of Satisfies Inclined Rail Field Geometry.
\end{proof}
\begin{definition}[Satisfies Terminal Rail Bar Laws]
  \label{def:physics:phyx-mini-0981:phyxminiproblems-problemphyxmini0981-satisfiesterminalrailbarlaws}
  \lean{PhyXMiniProblems.ProblemPhyXMini0981.SatisfiesTerminalRailBarLaws}
  \uses{def:physics:phyx-mini-0981:phyxminiproblems-problemphyxmini0981-lengthinmeters, def:physics:phyx-mini-0981:phyxminiproblems-problemphyxmini0981-massinkilograms, def:physics:phyx-mini-0981:phyxminiproblems-problemphyxmini0981-speedinmeterspersecond, def:physics:phyx-mini-0981:phyxminiproblems-problemphyxmini0981-accelerationmagnitudeinmeterspersecondsquared, def:physics:phyx-mini-0981:phyxminiproblems-problemphyxmini0981-accelerationinmeterspersecondsquared, def:physics:phyx-mini-0981:phyxminiproblems-problemphyxmini0981-magneticfluxdensityinteslas, def:physics:phyx-mini-0981:phyxminiproblems-problemphyxmini0981-electromotiveforceinvolts, def:physics:phyx-mini-0981:phyxminiproblems-problemphyxmini0981-electriccurrentinamperes, def:physics:phyx-mini-0981:phyxminiproblems-problemphyxmini0981-resistanceinohms, def:physics:phyx-mini-0981:phyxminiproblems-problemphyxmini0981-forceinnewtons, def:physics:phyx-mini-0981:phyxminiproblems-problemphyxmini0981-powerinwatts, def:physics:phyx-mini-0981:phyxminiproblems-problemphyxmini0981-inclinedrailbarsetup}
  School-level electromechanical laws at terminal speed. Each equation relates independent observables; none contains the requested closed power expression
\end{definition}
\begin{proof}
  This is the declared model definition of Satisfies Terminal Rail Bar Laws.
\end{proof}
% --- Archon declaration topology end ---
% --- Archon physics formalization source end ---
